\documentclass[10pt,a4paper]{amsart}
\usepackage[margin=19mm]{geometry}
\usepackage{amsmath,amssymb,mathtools}
\usepackage[scr=rsfs,cal=euler]{mathalfa}
\usepackage{xfrac}
\usepackage[unicode,hidelinks,bookmarks=false]{hyperref}
\newcommand{\ie}{i.e.\ }
\newcommand{\cf}{cf.\ }
\newcommand{\resp}{respectively\ }
\newcommand{\Subsection}{\subsection}
\providecommand{\nobreakdash}{\nobreak\hbox{-}}
\setlength{\emergencystretch}{2em}
\multlinegap0pt
\usepackage{amssymb}
\usepackage[all]{xy}
\usepackage{xcolor}
\usepackage{tikz}
\usepackage[shortlabels]{enumitem}
\newtheorem{lemma}{Lemma}
\newcommand{\C}{\mathbb{C}}
\newcommand{\Q}{\mathbb{Q}}
\newcommand{\R}{\mathbb{R}}
\newcommand\SL{\mathrm{SL}}
\newcommand{\Z}{\mathbb{Z}}
\newcommand{\pt}{\mathrm{pt}}
\title{Torus knots, the A-polynomial, and $\SL(2,\C)$}
\author{John A. Baldwin, Steven Sivek}
\date{Source: arXiv:2405.19197v2 (CC BY 4.0)}
\begin{document}
\maketitle

\noindent\emph{Source and licence.} Torus knots, the A-polynomial, and $\SL(2,\C)$. Version arXiv:2405.19197v2 (CC BY 4.0), distributed by arXiv under the Creative Commons Attribution 4.0 International licence, http://creativecommons.org/licenses/by/4.0/, as recorded in the arXiv licence metadata for this exact version and shown on the abstract page https://arxiv.org/abs/2405.19197v2. Full licence notice: \texttt{licenses/arxiv-2405.19197-CC-BY-4.0.txt}.\par\smallskip

\noindent\textit{Editorial scope.} Counted unit: the article's lemma lem:sl2-abelian-companion (source0.tex lines 887--889). The article's complete original proof of it is delivered with it (source0.tex lines 893--1013, one \texttt{proof} environment of 121 lines, which the article places away from the statement). No mathematics is written, repaired or re-derived: the statement and the proof are the article's own lines, in the article's own notation, and no retained line is substituted or re-flowed.  No further source-local context is needed: every cross-reference of every retained line resolves inside the two delivered spans.  Nothing was omitted from any retained span, so the package carries no omission record.  No retained line contains a citation, so the wrapper carries no bibliography and no bibliography entry.  No retained line contains a figure environment, an image-inclusion command or drawing code, so no image is delivered and the package declares no asset dependency.  Editorial formatting only, in addition: an AMS \texttt{amsart} preamble that restates the article's own declarations and macros used by the retained lines, namely the environments \texttt{lemma} on separate counters, together with the standard packages those lines need.  The retained environments print in this document as Lemma 1 in order of appearance, whereas the article prints the same items with the numbering of its own full document.  The complete native source of the article as distributed in the arXiv e-print for this exact version is bundled unchanged as \texttt{sources/159953/source0.tex}.
\begin{lemma} \label{lem:sl2-abelian-companion}
Let $K = P(C)$ be a satellite knot whose pattern $P$ has winding number $w \geq 1$, and let $r\in\Q$.  If $r$-surgery on $C$ is not $\SL(2,\C)$-abelian, then neither is $rw^2$-surgery on $K$.
\end{lemma}

\begin{proof}[Proof of Lemma \ref{lem:sl2-abelian-companion}]
Write $r = p/q$, with $p$ and $q$ relatively prime and $q > 0$.  Let $d=\gcd(q,w^2)$, and note that
\[ rw^2 = \frac{pw^2/d}{q/d} \]
in lowest terms, since $\gcd(pw^2/d, q/d) = 1$.  Let $E_C = S^3 \setminus \nu(C)$ denote the exterior of $C$, and fix a nonabelian representation
\begin{align*}
\rho: \pi_1(E_C) &\to \SL(2,\C) \\
\mu_C^p \lambda_C^q &\mapsto 1,
\end{align*}
which exists by hypothesis.  We will assume, by replacing $\rho$ with a conjugate if necessary, that $\rho(\mu_C)$ is in Jordan normal form.  This means that $\rho(\mu_C)$ is a diagonal matrix unless its eigenvalues are equal, in which case these equal eigenvalues are  both $+1$ or both $-1$  since $\det(\rho(\mu_C)) = 1$.  We also note that $\rho(\mu_C)$ cannot be either of the central elements $\pm I$, since $\rho(\mu_C)$  normally generates the nonabelian image of $\rho$.

Letting $V = (S^1\times D^2) \setminus \nu(P)$, we may decompose the exterior $E_K = S^3\setminus \nu(K)$ as
\[ E_K \cong E_C \cup_T V. \]
Our goal will be to extend $\rho$ across $\pi_1(V)$, by constructing an \emph{abelian} representation
\[ \rho_V: \pi_1(V) \to \SL(2,\C) \]
such that
\begin{align*}
\rho_V(\mu_C) = \rho_V( [\{\pt\} \times \partial D^2] ) &= \rho(\mu_C) \\
\rho_V(\lambda_C) = \rho_V( [S^1 \times \{\pt\}] ) &= \rho(\lambda_C) \\
\rho_V(\mu_P^{pw^2/d} \lambda_P^{q/d}) &= 1.
\end{align*}
Then $\rho$ and $\rho_V$ will glue together to give a representation \[\pi_1(E_K) \to \SL(2,\C)\] that descends to a nonabelian representation of $\pi_1(S^3_{rw^2}(K))$, as desired.

Observe that any abelian representation $\rho_V: \pi_1(V) \to \SL(2,\C)$ must factor through 
\[ H_1(V;\Z) \cong \Z^2, \]
which is generated by the classes $[\mu_P]$ and $[\lambda_C]$ subject to the relations
\[
[\mu_C] = w [\mu_P]\, \textrm{ and }\,
[\lambda_P] = w[\lambda_C].
\]
The desired $\rho_V$ is therefore determined by the elements $\rho_V(\mu_P)$ and $\rho_V(\lambda_C)$, and the value $\rho_V(\lambda_C) = \rho(\lambda_C)$ is already fixed by $\rho$ -- we note that this implies that $\rho_V(\lambda_P) = \rho(\lambda_C)^w$ -- so we wish to find $\rho_V(\mu_P) \in \SL(2,\C)$ satisfying
\begin{align*}
\rho_V(\mu_P)^w &= \rho(\mu_C) \\
\rho_V(\mu_P^{pw^2/d} \lambda_P^{q/d}) &= 1.
\end{align*}
We consider several cases separately, recalling that $\rho(\mu_C)$ is in Jordan normal form.

\vspace{1em}\noindent
\underline{Case 1}: $\rho(\mu_C) \in \SL(2,\C)$ is a diagonal matrix $\left(\begin{smallmatrix} \alpha & 0 \\ 0 & \alpha^{-1} \end{smallmatrix}\right)$, where $\alpha \neq \pm1$.\\

In this case, $\rho(\mu_C)$ is diagonal but not central, and $\rho(\lambda_C)$ commutes with it, so $\rho(\lambda_C)$ must be diagonal as well: we can write
\begin{align*}
\rho(\mu_C) &= \begin{pmatrix} \alpha & 0 \\ 0 & \alpha^{-1} \end{pmatrix}, &
\rho(\lambda_C) &= \begin{pmatrix} \beta & 0 \\ 0 & \beta^{-1} \end{pmatrix}
\end{align*}
where $\alpha^p \beta^q = 1$.  We let $\alpha = s e^{i\theta}$ and $\beta = t e^{i\phi}$, where $s$ and $t$ are positive real numbers and $\theta,\phi\in\R$, and then the condition $\alpha^p\beta^q=1$ means that
\[
s^pt^q = 1 \,\textrm{ and }\,
p\theta + q\phi = 2\pi m
\]
for some integer $m$.  Let us write
\[ \eta = s^{1/w} \exp\left( \frac{(\theta+2\pi k)i}{w} \right) \]
for some integer $k$, chosen so that 
\[ m+pk \equiv 0 \pmod{d}; \]
we note that this is possible since $d=\gcd(q,w^2)$ divides $q$ and is thus coprime to $p$, meaning that $p$ is invertible $\!\!\!\pmod{d}$.   Then $\eta^w = \alpha$, so we set
\begin{align*}
\rho_V(\mu_P) &= \begin{pmatrix} \eta & 0 \\ 0 & \eta^{-1} \end{pmatrix}, &
\rho_V(\lambda_P) &= \begin{pmatrix} \beta^w & 0 \\ 0 & \beta^{-w} \end{pmatrix}
\end{align*}
so that $\rho_V(\mu_P)^w = \rho(\mu_C)$ and $\rho_V(\lambda_P) = \rho(\lambda_C)^w$.

Having done this, we now compute that
\[ \rho_V(\mu_P^{pw^2/d} \lambda_P^{q/d}) = \begin{pmatrix} z & 0 \\ 0 & z^{-1} \end{pmatrix}, \]
where $z = \eta^{pw^2/d} \cdot (\beta^w)^{q/d}$ satisfies
\begin{align*}
z &= s^{pw/d} t^{qw/d} \cdot \exp\left( \frac{pw(\theta+2\pi k)i}{d} + \frac{\phi i \cdot qw}{d}\right) \\
&= (s^pt^q)^{w/d} \cdot \exp\left( \frac{\big((p\theta+q\phi) + 2\pi pk\big)wi }{d} \right) \\
&= \exp\left( 2\pi \left(\frac{m+pk}{d}\right) \cdot wi \right)
\end{align*}
since $s^pt^q = 1$ and $p\theta + q\phi = 2\pi m$.  But we chose $k$ so that $(m+pk)/d$ would be an integer, so $z=1$ and thus $\rho_V(\mu_P^{pw^2/d}\lambda_P^{q/d}) = 1$, meaning that $\rho_V$ is the desired representation.

\vspace{1em}\noindent
\underline{Case 2}: $\rho(\mu_C)$ is a $2\times 2$ Jordan block with  eigenvalues which are either both $1$ or both $-1$, where  in the latter case we require that $w$ is odd.\\

Since $\rho(\mu_C) \neq \pm1$ commutes with $\rho(\lambda_C)$,  the latter is also a Jordan block. So, we have
\begin{align*}
\rho(\mu_C) &= \epsilon \begin{pmatrix} 1 & a \\ 0 & 1 \end{pmatrix}, &
\rho(\lambda_C) &= \eta \begin{pmatrix} 1 & b \\ 0 & 1 \end{pmatrix}
\end{align*}
for some $a,b\in\C$ with $a\neq 0$ and some $\epsilon,\eta=\pm1$ with $\epsilon^w = \epsilon$.  Then $\rho(\mu_C^p\lambda_C^q)=1$ is equivalent to $\epsilon^p\eta^q=1$ and $ap+bq=0$.  We define $\rho_V$ by
\begin{align*}
\rho_V(\mu_P) &= \epsilon \begin{pmatrix} 1 & a/w \\ 0 & 1 \end{pmatrix}, &
\rho_V(\lambda_P) &= \eta^w \begin{pmatrix} 1 & bw \\ 0 & 1 \end{pmatrix},
\end{align*}
which satisfies \[\rho_V(\mu_P)^w = \rho(\mu_C) \, \textrm{ and }\,\rho_V(\lambda_P) = \rho(\lambda_C)^w,\] and we compute that
\begin{align*}
\rho_V(\mu_P^{pw^2/d}\lambda_P^{q/d}) &= \epsilon^{pw^2/d} \begin{pmatrix} 1 & apw/d \\ 0 & 1 \end{pmatrix} \cdot \eta^{w(q/d)} \begin{pmatrix} 1 & bqw/d \\ 0 & 1 \end{pmatrix} \\
&= \epsilon^{pw^2/d}\eta^{w(q/d)} \begin{pmatrix} 1 & (ap+bq)w/d \\ 0 & 1 \end{pmatrix} \\
&= \epsilon^{pw^2/d}\eta^{w(q/d)} \begin{pmatrix} 1 & 0 \\ 0 & 1 \end{pmatrix}
\end{align*}
since $ap+bq=0$.  If $d$ is odd, then the leading coefficient on the right satisfies
\[ \epsilon^{pw^2/d}\eta^{w(q/d)} = \left(\epsilon^{pw^2/d}\eta^{w(q/d)}\right)^d = (\epsilon^w)^{pw} \eta^{qw} = \epsilon^{pw}\eta^{qw} = (\epsilon^p\eta^q)^w = 1. \]
Otherwise $d=\gcd(q,w^2)$ is even, hence both $q$ and $w$ must be even; but then $\eta^q = \eta^w=1$, and $\epsilon^p\eta^q=1$ implies that $\epsilon^p=1$, so that
\[ \epsilon^{pw^2/d}\eta^{w(q/d)} = (\epsilon^p)^{w^2/d} (\eta^w)^{q/d} = 1. \]
In either case we have $\rho_V(\mu_P^{pw^2/d}\lambda_P^{q/d}) = 1$, so $\rho_V$ is the desired representation.

\vspace{1em}\noindent
\underline{Case 3}: $\rho(\mu_C)$ is a $2\times 2$ Jordan block with eigenvalues $-1$, and $w$ is even.\\

As in the previous case, we can write
\begin{align*}
\rho(\mu_C) &= -\begin{pmatrix} 1 & a \\ 0 & 1 \end{pmatrix}, &
\rho(\lambda_C) &= \eta \begin{pmatrix} 1 & b \\ 0 & 1 \end{pmatrix}
\end{align*}
for some $a,b\in\C$ with $a\neq 0$ and some $\eta=\pm1$, and then $\rho(\mu_C^p\lambda_C^q)=1$ is equivalent to $(-1)^p\eta^q=1$ and $ap+bq=0$.  The problem is that now the matrix $\rho(\mu_C)$ does not have a $w$th root.  We fix this by considering the nontrivial central character
\[ \chi: \pi_1(E_C) \twoheadrightarrow H_1(E_C) \to \{\pm1\} \]
that sends $\mu_C$ to $-1$ and $\lambda_C$ to $+1$, and setting $\rho' = \chi\cdot\rho$, so that
\begin{align*}
\rho'(\mu_C) &= \begin{pmatrix} 1 & a \\ 0 & 1 \end{pmatrix}, &
\rho'(\lambda_C) &= \eta \begin{pmatrix} 1 & b \\ 0 & 1 \end{pmatrix}.
\end{align*}
We can extend $\rho'$, rather than the original $\rho$, across $\pi_1(V)$ by setting
\begin{align*}
\rho_V(\mu_P) &= \begin{pmatrix} 1 & a/w \\ 0 & 1 \end{pmatrix}, &
\rho_V(\lambda_P) &= \begin{pmatrix} 1 & bw \\ 0 & 1 \end{pmatrix},
\end{align*}
and we have $\rho_V(\mu_P)^w = \rho'(\mu_C)$ and $\rho_V(\lambda_P) = \rho'(\lambda_C)^w$ since $w$ is even.  Now 
\[ \rho_V(\mu_P^{pw^2/d}\lambda_P^{q/d}) = \begin{pmatrix} 1 & (ap+bq)w/d \\ 0 & 1 \end{pmatrix} = \begin{pmatrix} 1 & 0 \\ 0 & 1 \end{pmatrix}, \]
so $\rho_V$ is the desired extension of $\rho'$, and it has nonabelian image since $\rho'$ does.

\vspace{1em}These three cases cover all possible values of $\rho(\mu_C)$, so we conclude that $rw^2$-surgery on $K$ cannot be $\SL(2,\C)$-abelian after all.
\end{proof}

\end{document}
